\par\medskip\noindent\textbf{Cómo leer la salida.}
{\small\setlength{\tabcolsep}{4pt}\renewcommand{\arraystretch}{1.15}
\begin{xltabular}{\linewidth}{@{}>{\raggedright\arraybackslash}p{0.24\linewidth}>{\raggedleft\arraybackslash}p{0.15\linewidth}>{\raggedright\arraybackslash}X@{}}
\toprule Indicador & Valor & Interpretación \\ \midrule \endfirsthead
\toprule Indicador & Valor & Interpretación \\ \midrule \endhead
\bottomrule \endlastfoot
\texttt{alfa anual} & 6,78 a 21,79\,\% & Rendimiento anual no explicado por el mercado. \\
\texttt{t(alfa)} & 0,89 a 1,96 & Con Newey-West ninguna alfa individual es claramente significativa (Amazon en el límite): cinco años son pocos. \\
\texttt{beta} & 1,13 a 1,34 & Las cuatro acciones amplifican los movimientos del mercado. \\
\texttt{t(beta = 1)} & 2,88 a 4,79 & Todas las betas son significativamente mayores que 1. \\
\texttt{R2} & 0,31 a 0,45 & El mercado explica entre un tercio y la mitad de la varianza diaria de cada acción. \\
\texttt{vol. idiosincrática} & 17,4 a 25,4\,\% & Volatilidad anual propia de cada acción, diversificable. \\
Cartera 1/N & alfa 13,38\,\% ($t=2{,}23$); beta 1,229 & Al diversificar, el riesgo propio baja y la alfa de la cartera se vuelve significativa. \\
\texttt{R2} de la cartera & 0,566 & Mayor que el de cada acción: la diversificación elimina parte del riesgo idiosincrático. \\
Promedio de las betas & 1,2286 & Igual a la beta de la cartera: la beta de una cartera es el promedio ponderado de las betas. \\
Estabilidad & 1,185 / 1,261 & La beta subió ligeramente entre 2014--2015 y 2016--2018. \\
\end{xltabular}}
